\documentclass[11pt]{amsart}

\usepackage[T1]{fontenc}
\usepackage{lmodern}
\usepackage{microtype}
\usepackage{amsmath,amssymb,amsthm}
\usepackage[hidelinks]{hyperref}

\hypersetup{
  pdftitle={A Counterexample to Marin's Conjecture 5.14}
}

\newtheorem{theorem}{Theorem}
\theoremstyle{remark}
\newtheorem{remark}[theorem]{Remark}

\newcommand{\ctop}{c_{\mathrm{top}}}
\newcommand{\Ptwo}{\mathcal P^{(2)}}

\title[Counterexample to Mar\'\i n's Conjecture 5.14]
{A Counterexample to Mar\'\i n's Conjecture 5.14}

\date{}

\subjclass[2020]{05E10, 20G05}
\keywords{
Littlewood restriction, orthogonal branching, Schur polynomials,
counterexample
}

\begin{document}

\begin{abstract}
Conjecture~5.14 of Mar\'\i n, in arXiv:2608.18302v1, asserts that when
the maximal top weight of a certain odd-orthogonal branching expansion is
attained by more than one constituent, the tied top contributions cancel.
At $(t,r)=(3,2)$, we show that
\[
\beta=(12,11,10,3,2,1,0)
\]
has exactly three maximizing constituents with nonzero orthogonal
restriction coefficient, namely $(6,6,6)$, $(6,6,2)$, and $(6,6,0)$.
Each restriction coefficient is $1$ and each top coefficient is
$\pm1$, so the tied sum is odd and cannot vanish. This refutes the
conjecture.
\end{abstract}

\maketitle

\section{The counterexample}

We follow the notation of Mar\'\i n~\cite{Marin}. For a constituent
$\Lambda$ with orthogonal restriction coefficient $a^B_\Lambda\ne0$, let
$T(\Lambda)$ be ``the top weight of $c(\Lambda,\cdot)$ delivered by
Proposition~5.10(v)'' and let
\[
\ctop(\Lambda)=c(\Lambda,T(\Lambda))
\]
be ``its value there''. Then
\[
M(\beta)=\max\{T(\Lambda):a^B_\Lambda\ne0\}
\]
is ``the largest of them, the maximum taken in the order used by''
Mar\'\i n's formula for $\mu_{\max}$. Conjecture~5.14 of~\cite{Marin} reads:
if $M(\beta)$ is attained by more than one $\Lambda$, then
\begin{equation}\label{eq:conjecture}
\sum_{\Lambda\,:\,T(\Lambda)=M(\beta)}
a^B_\Lambda\ctop(\Lambda)=0.
\end{equation}

Four clauses of~\cite{Marin} are used below, and we quote them so that the
proof can be followed without the source. Write $t=2m'+1$. The entries of
$V=2(\Lambda+\rho_{B_{R'}})$ are sorted into folded residue classes modulo
$t$, of which there are $m'$ nonzero ones and one class $0$, and
$n_j=\#\{i:V_i\text{ folds to class }j\}$; a set $S$ of indices is a
\emph{transversal} when it ``contains exactly one index from each nonzero
class and none from class $0$''. Proposition~5.10 of~\cite{Marin}
describes the numerator $\nu(\Lambda,\cdot)$ attached to $\Lambda$ in these
terms: (iii) ``$\nu\equiv0$ if and only if some $n_j=0$ --- some nonzero
residue class is missed by $V=2(\Lambda+\rho_{B_{R'}})$''; (iv) ``the value
is explicit: at the point indexed by $S$ and a chirality it is
$\epsilon_t\,\det(w)\,\delta_S\in\{\pm1\}$'', the three factors being those
of the source, of which we use only the membership in $\{\pm1\}$; and (v)
``Let $S_{\min}$ be the transversal taking, in each nonzero class, the
index with the \emph{largest} label --- equivalently the smallest value of
$V$. Then the alternant index of $(S_{\min},+)$ \emph{dominates
coordinatewise} the index of every other point of
$\mathrm{supp}\,\nu(\Lambda,\cdot)$, strictly in at least one coordinate.
In particular it is the unique maximiser of $\langle w_0,\cdot\rangle$ for
\emph{every} $w_0$ with all coordinates positive.''

The indexing is the one fixed in the proof of that clause: alternants are
indexed ``by the doubled vector $2(\mu+\rho_{D_r})$, so that the point of
$\nu$ attached to $(S,\text{chirality})$ has index $V|_{S^{\mathrm c}}$
sorted decreasingly''. At $t=3$ one has $m'=1$, so there is a single
nonzero folded class, a transversal is a single index, and $S_{\min}$ is
the index carrying the smallest entry of $V$ in that class. The top weight
attached to $\Lambda$ is thus read off the vector obtained from $V(\Lambda)$
by deleting that entry, which is the rule used below.

The passage from the numerator to $c$ is a division by the specialised
denominator $\Delta_t$, and Proposition~5.19(ii) of~\cite{Marin} inverts it
in closed form: with alternants indexed as above and $\tilde\nu$,
$\tilde c$ the $W(D_r)$-antisymmetric extensions of $\nu$ and $c$,
multiplication by $\Delta_t$ is
``$a_X\mapsto\sum_{\varepsilon\in\{\pm1\}^r}
(\prod_j\varepsilon_j)\,a_{X+t\varepsilon}$, and it is inverted by''
\[
\tilde c(X)=s_r\!\!\sum_{k\in(2\mathbb{Z}_{\ge0}+1)^r}\!\!
\tilde\nu\bigl(X+tk\bigr),
\]
``a sum with finitely many nonzero terms'', $s_r$ being a global sign. The
shifts $tk$ over odd $k$ are the ones compared in the proof below.

\begin{theorem}
For $(t,r)=(3,2)$ and
\[
\beta=(12,11,10,3,2,1,0),
\]
one has $M(\beta)=(6,6)$. The constituents with
$a^B_\Lambda\ne0$ and $T(\Lambda)=M(\beta)$ are exactly
\[
(6,6,6),\qquad (6,6,2),\qquad (6,6,0),
\]
and each has $a^B_\Lambda=1$. Consequently, the sum in
\eqref{eq:conjecture} is nonzero.
\end{theorem}

\begin{proof}
Here $N=7$, $\delta=(6,5,4,3,2,1,0)$, and
\[
\lambda=\beta-\delta=(6,6,6).
\]
Since $\ell(\lambda)=3\le 7/2$ we are in the stable range, where
Littlewood's orthogonal restriction formula gives
\begin{equation}\label{eq:littlewood}
a^B_\Lambda
=
\sum_{\gamma\in\Ptwo}
c^{(6,6,6)}_{\gamma,\Lambda},
\end{equation}
where $\Ptwo$ is the set of partitions with even parts. Outside the stable
range this formula is replaced by a finer branching rule, for which
see~\cite{JangKwon}; the shape at hand does not need it.

For a partition
$\Lambda=(\Lambda_1,\Lambda_2,\Lambda_3)$ contained in the rectangle
$(6,6,6)$, the $GL_3$-identity
$V_{(6,6,6)}=\det^6$ gives the rectangular duality formula
\[
c^{(6,6,6)}_{\gamma,\Lambda}
=
\begin{cases}
1,
&
\gamma=(6-\Lambda_3,6-\Lambda_2,6-\Lambda_1),
\\
0,
&
\text{otherwise}.
\end{cases}
\]
Since $c^{(6,6,6)}_{\gamma,\Lambda}=0$ unless $\Lambda$ is contained in
$(6,6,6)$ --- and $\Lambda$ has at most $R'=3$ rows in any case, being a
dominant $B_3$ weight --- no $\Lambda$ outside that rectangle contributes to
\eqref{eq:littlewood}. Thus
\begin{equation}\label{eq:branching}
a^B_\Lambda
=
\begin{cases}
1,
&
6\ge\Lambda_1\ge\Lambda_2\ge\Lambda_3\ge0
\text{ and all }\Lambda_i\text{ are even},
\\
0,
&
\text{otherwise}.
\end{cases}
\end{equation}

For every $\Lambda$ occurring in \eqref{eq:branching}, set
\[
V(\Lambda)
=
2(\Lambda+\rho_{B_3})
=
(2\Lambda_1+5,2\Lambda_2+3,2\Lambda_3+1).
\]
At $t=3$ there is one nonzero folded residue class, represented by
$\{\pm1\pmod3\}$. By Proposition~5.10(iii) of~\cite{Marin}, the
numerator attached to $\Lambda$ vanishes only if all three entries of
$V(\Lambda)$ are divisible by $3$. This cannot occur here. Indeed,
if all three entries were divisible by $3$, then
\[
\Lambda_1\equiv2,\qquad
\Lambda_2\equiv0,\qquad
\Lambda_3\equiv1
\pmod3.
\]
Together with evenness and
$0\le\Lambda_3\le\Lambda_1\le6$, this would force
$\Lambda_1=2$ and $\Lambda_3=4$, a contradiction.

Let $X(\Lambda)$ be the decreasing two-vector obtained by deleting
from $V(\Lambda)$ the smallest entry lying in the nonzero folded
class. Propositions~5.10(iv),(v) and~5.19 of~\cite{Marin} imply
\begin{equation}\label{eq:top}
2\bigl(T(\Lambda)+\rho_{D_2}\bigr)
=
X(\Lambda)-3(1,1),
\qquad
\ctop(\Lambda)\in\{\pm1\},
\end{equation}
where $\rho_{D_2}=(1,0)$.

Indeed, Proposition~5.10(iv),(v) gives the unique coordinatewise
maximal numerator index $X(\Lambda)$ with coefficient $\pm1$. In the
inversion formula of Proposition~5.19(ii), the shift $(1,1)$ contributes
this index at $X(\Lambda)-3(1,1)$. Every other shift, after
$D_2$-straightening, strictly dominates $X(\Lambda)$ coordinatewise
and therefore misses the numerator support. The same argument excludes
all quotient indices above $X(\Lambda)-3(1,1)$.

Since $\Lambda\subseteq(6,6,6)$, the entries of $V(\Lambda)$ are
strictly decreasing and bounded by $(17,15,13)$. Hence
\[
X(\Lambda)\le(17,15)
\]
coordinatewise, and \eqref{eq:top} gives
\[
T(\Lambda)\le(6,6).
\]
Equality requires $X(\Lambda)=(17,15)$, which forces
$\Lambda_1=\Lambda_2=6$. Write $\Lambda=(6,6,k)$. By
\eqref{eq:branching}, its restriction coefficient is nonzero exactly
for
\[
k\in\{0,2,4,6\}.
\]

Now
\[
V(6,6,k)=(17,15,2k+1).
\]
For $k=0,2,6$, the third entry lies in the nonzero folded class and
is smaller than $17$, so it is deleted and
\[
X=(17,15).
\]
For $k=4$, the third entry is $9\equiv0\pmod3$; the deleted entry is
then $17$, and
\[
X=(15,9).
\]
Therefore the three and only three maximizers with nonzero restriction
coefficient are
\[
(6,6,6),\qquad (6,6,2),\qquad (6,6,0).
\]
Their restriction coefficients are all $1$, while their top
coefficients lie in $\{\pm1\}$ by \eqref{eq:top}.

The comparisons just made are coordinatewise, whereas the maximum defining
$M(\beta)$ is taken in the Weyl majorisation order
$\mu\preceq\mu'\iff\mu\in\mathrm{conv}\bigl(W(C_r)\,\mu'\bigr)$. At $r=2$
this costs nothing: $\mathrm{conv}\bigl(W(C_2)(6,6)\bigr)=[-6,6]^2$, and the
entries of $V(\Lambda)$, hence of $X(\Lambda)$, are positive, so
\eqref{eq:top} together with $T(\Lambda)\le(6,6)$ puts every $T(\Lambda)$
in $[-2,6]^2$. Thus
$T(\Lambda)\preceq(6,6)$ for every $\Lambda$, and $(6,6)$ is the maximum in
that order as well; being the coordinatewise maximum, it also maximises
every linear functional with positive coefficients.

The left-hand side of \eqref{eq:conjecture} is therefore a sum of three
numbers in $\{\pm1\}$ and cannot be zero.
\end{proof}

\begin{remark}[Scope]
The hypothesis of Conjecture~5.14 is that $M(\beta)$ be ``attained by more
than one $\Lambda$'', and $\beta=(12,11,10,3,2,1,0)$ satisfies it. The
discussion around the conjecture sometimes phrases the tie case as a top
attained twice: Remark~5.13 of~\cite{Marin} writes ``Where it is attained
twice, the question is whether the tie always cancels''. The dichotomy that
remark draws is nevertheless between a top attained once and a top not
attained once. Its biconditional, over the $132$ shapes tested there, is
that the prediction $M(\beta)$ for the top weight of the composite is
``correct \emph{exactly} on the $105$ where the largest top is attained by
a single $\Lambda$''; and the map of conjectures in~\cite{Marin} records
that Conjecture~5.14, with Proposition~5.10(v), gives (L2) ``\emph{only}
where the top is attained once''. A threefold tie is a shape where the top
is not attained once, which is where the conjecture is called upon, so the
example refutes it as stated.
\end{remark}

\begin{remark}[The verified population]
Conjecture~5.14 is reported in~\cite{Marin} as machine-checked, and
$\beta=(12,11,10,3,2,1,0)$ lies outside the population checked. That
population is the one of Remark~5.13, ``the $132$ shapes at
$(t,r)=(3,2),(5,2),(3,3)$''; the verification table records
Conjecture~5.14 there as ``$27/27$ tied shapes; the biconditional
$132/132$''. Those shapes are the $123$
nonvanishing forms of Observation~5.27 --- $82$ at $(3,2)$ and $41$ at
$(5,2)$ --- together with the $9$ at $(3,3)$, and the box of
Observation~5.27 is ``shapes with $\max\beta\le9$ at $(t,r)=(3,2)$ and
$\max\beta\le10$ at $(5,2)$''. Here $\max\beta=12$, so the example
contradicts no computation reported in~\cite{Marin}.
\end{remark}

\begin{thebibliography}{9}

\bibitem{JangKwon}
I.-S. Jang and J.-H. Kwon,
\emph{Flagged Littlewood--Richardson tableaux and branching rule for
classical groups},
J. Combin. Theory Ser. A \textbf{181} (2021),
Paper No. 105419,
\href{https://doi.org/10.1016/j.jcta.2021.105419}
{doi:10.1016/j.jcta.2021.105419}.

\bibitem{Marin}
C. Mar\'\i n,
\emph{Schur polynomials twisted by roots of unity and reciprocal pairs:
torsion filters, fusion quotients, and an equal-rank reduction at odd
order},
\href{https://arxiv.org/abs/2608.18302v1}
{arXiv:2608.18302v1 [math.CO]} (2026).

\end{thebibliography}

\end{document}